\expandafter\ifx\csname CAPETINCLUDEMODE\endcsname\relax
\documentclass[11pt,a4paper]{article}
% preamble.tex — packages et commandes communes a toutes les fiches CAPET SII IE
\input{glyphtounicode}
\pdfgentounicode=1
\usepackage{cmap}
\usepackage[T1]{fontenc}
\usepackage[utf8]{inputenc}
\usepackage[french]{babel}
\usepackage[a4paper, top=1.6cm, bottom=1.6cm, left=1.4cm, right=1.4cm]{geometry}
\setlength{\headheight}{14.5pt}
\usepackage{amsmath, amssymb, amsfonts}
\usepackage{siunitx}
\sisetup{locale=FR, output-decimal-marker={,}, inter-unit-product=\ensuremath{{}\cdot{}}}
\usepackage{array, booktabs, multirow, tabularx}
\usepackage{xcolor}
\usepackage[most, breakable]{tcolorbox}
\usepackage{multicol}
\usepackage{enumitem}
\usepackage{mathtools}
\usepackage{lmodern}
\usepackage{microtype}
\usepackage{fancyhdr}
\usepackage{tikz}
\usetikzlibrary{positioning,arrows.meta}
\usepackage{lastpage}

% ── Couleurs par domaine ──────────────────────────────────────────────
\definecolor{colElec}{RGB}{26, 82, 148}       % Électrotechnique — bleu foncé
\definecolor{colElecLight}{RGB}{210, 225, 245} % fond bleu clair
\definecolor{colAuto}{RGB}{34, 120, 66}        % Automatique — vert
\definecolor{colMeca}{RGB}{185, 90, 20}        % Mécanique — orange
\definecolor{colTherm}{RGB}{160, 40, 40}       % Thermique — rouge
\definecolor{colQT}{RGB}{232, 244, 233}        % fond question type (vert très léger)
\definecolor{colJury}{RGB}{255, 243, 205}      % fond jury (jaune ambre)
\definecolor{colJuryFrame}{RGB}{200, 130, 0}   % bordure jury
\definecolor{colGrayLight}{RGB}{245, 245, 247} % fond gris neutre
\definecolor{colSeparator}{RGB}{180, 180, 180} % gris séparateur

% ── Style tcolorbox de section principale ────────────────────────────
\tcbset{
  ficheSection/.style={
    breakable,
    enhanced,
    colback=#1!7!white,
    colframe=#1,
    colbacktitle=#1,
    coltitle=white,
    fonttitle=\bfseries\normalsize,
    left=4pt, right=4pt, top=3pt, bottom=4pt,
    before skip=6pt, after skip=6pt,
  }
}

% ── Style encadré question type ──────────────────────────────────────
\tcbset{
  questiontype/.style={
    breakable,
    enhanced,
    colback=colQT,
    colframe=colElec!55!black,
    colbacktitle=colElec!55!black,
    coltitle=white,
    fonttitle=\bfseries\small,
    left=4pt, right=4pt, top=3pt, bottom=3pt,
    before skip=5pt, after skip=5pt,
  }
}

% ── Style encadré jury ───────────────────────────────────────────────
\tcbset{
  jury/.style={
    breakable,
    enhanced,
    colback=colJury,
    colframe=colJuryFrame,
    fonttitle=\bfseries\normalsize,
    left=4pt, right=4pt, top=3pt, bottom=3pt,
  }
}

% ── Mise en page ─────────────────────────────────────────────────────
\pagestyle{fancy}
\fancyhf{}
\renewcommand{\headrulewidth}{0.4pt}
\renewcommand{\footrulewidth}{0.3pt}

\setlist[itemize]{leftmargin=1.2em, itemsep=1pt, topsep=2pt, parsep=0pt}
\setlist[enumerate]{leftmargin=1.2em, itemsep=3pt, topsep=2pt, parsep=0pt}

\setlength{\parskip}{2pt}
\setlength{\columnsep}{10pt}

\def\CAPETBEGINFILE{\begin{document}}
\def\CAPETENDFILE{\end{document}}
\else
\def\CAPETBEGINFILE{}
\def\CAPETENDFILE{}
\fi

\fancyhead[L]{\textbf{\textcolor{colElec}{T24}} - Systemes embarques - Microcontreleurs, API, FPGA}
\fancyhead[R]{\small\textit{CAPET SII IE - Session 2026}}
\fancyfoot[C]{\small\thepage\,/\,\pageref{LastPage}}

\CAPETBEGINFILE

% -----------------------------------------------------------------------
%  BLOC EN-TeTE
% -----------------------------------------------------------------------

\begin{tcolorbox}[
  enhanced, breakable,
  colback=colElec,
  colframe=colElec,
  left=6pt, right=6pt, top=5pt, bottom=5pt,
  before skip=0pt, after skip=4pt
]
  {\color{white}\bfseries\large T24 \quad Systemes embarques - Microcontreleurs, API, FPGA}\\[2pt]
  {\color{colElec!30!white}\small
    \textbf{Domaine :} electronique - Commande numerique et systemes programmables \quad
    \textbf{Presence :} \textbf{8/13} sujets (2013-2026) - theme \emph{tres frequent}
  }\\[1pt]
  {\color{colElec!20!white}\footnotesize
    Sous-themes : microcontreleur (Arduino, PIC) - registre ADCON0 - API S7-1500 - PWM - bus CAN/OneWire/LoRa - DSP - FPGA/VHDL (programme) - mise - l'echelle automate
  }
\end{tcolorbox}

% -----------------------------------------------------------------------
%  1. RAPPELS DE COURS
% -----------------------------------------------------------------------

\begin{tcolorbox}[
  ficheSection=colElec,
  title={1. Rappels de cours},
]

% -- 1.1 Architecture systemes embarques -------------------------------
\textbf{1.1 Architecture des systemes embarques}

\smallskip
\begin{center}
\renewcommand{\arraystretch}{1.25}
\begin{tabular}{llll}
\toprule
\textbf{Composant} & \textbf{Rele} & \textbf{Exemples CAPET} & \textbf{Prog.} \\
\midrule
Microcontreleur & CPU + peripheriques integres & Arduino, PIC 18F8723 & C/Python \\
DSP & Traitement signal haute vitesse & DSP 40 MHz (2024) & C optimise \\
API/PLC & Contrele industriel robuste & Siemens S7-1500 (2024) & Ladder/Grafcet \\
FPGA & Circuit reconfigurable parallele & (Programme SII) & VHDL \\
\bottomrule
\end{tabular}
\end{center}

\medskip\hrule\vspace{4pt}

% -- 1.2 Microcontreleur - CAN interne --------------------------------
\textbf{1.2 Microcontreleur - CAN interne (PIC 18F8723)}

\begin{itemize}
  \item 1 CAN partage entre 16 entrees analogiques AN0eAN15 via multiplexeur
  \item Selection de la voie : registre ADCON0, champs CHS3:CHS0
  \item Resolution 10 bits ; Quantum $q = V_{ref}/2^{10} = V_{ref}/1024$
  \item Demarrage conversion : bit GO/DONE = 1
\end{itemize}

\medskip\hrule\vspace{4pt}

% -- 1.3 PWM ----------------------------------------------------------
\textbf{1.3 PWM - Modulation de Largeur d'Impulsion}

Signal creneau - frequence fixe (50 Hz pour servomoteur), rapport cyclique $\alpha = t_{on}/T$.

\[
  t_{on} = \frac{\theta - \theta_{min}}{\theta_{max} - \theta_{min}}\,(t_{on,max} - t_{on,min}) + t_{on,min}
\]

Frequence d'echantillonnage pour resolution angulaire $\delta\theta$ :
\[
  f_{ech} = \frac{\theta_{max} - \theta_{min}}{\delta\theta \cdot T} = \frac{N_{positions}}{T}
  \quad \Rightarrow \quad
  f_{ech} = N_{pos} \times f_{PWM}
\]
Exemple (2025) : $\delta\theta = 0{,}1^\circ$, plage $180^\circ$ - $N = 1800$ - $f_{ech} = 1{,}8\,\mathrm{MHz}$ (PWM - 50 Hz).

\medskip\hrule\vspace{4pt}

% -- 1.4 API Siemens S7-1500 ------------------------------------------
\textbf{1.4 Automate S7-1500 - Mise - l'echelle}

\[
  T = T_{min} + \frac{PIW8 - 0}{27\,648 - 0}\,(T_{max} - T_{min})
\]
\begin{itemize}
  \item Module AI : PIW8 $\in [0\,;\,27\,648]$ pour $I_s \in [4\,;\,20]\,\mathrm{mA}$
  \item Programme cyclique : lecture entrees - traitement - ecriture sorties
  \item Capteur OneWire DS1820 : commande 0x44 (conversion) ; 0xBE (lecture scratchpad)
\end{itemize}

\medskip\hrule\vspace{4pt}

% -- 1.5 DSP ----------------------------------------------------------
\textbf{1.5 DSP - Calcul de frequence maximale}

\[
  f_{max} = \frac{P \cdot n}{60}
  \quad P : \text{points/tour codeur}, \quad n : \text{vitesse (tr/min)}
\]
Condition : $f_{DSP} \geq f_{max}$. Exemple (2024) : $f_{max} = 1326\,\mathrm{Hz}$, $f_{DSP} = 40\,\mathrm{MHz}$ - .

\medskip\hrule\vspace{4pt}

% -- 1.6 Bus CAN Controller Area Network -------------------------------
\textbf{1.6 Bus CAN (Controller Area Network)}

Bus serie differentiel multi-maetres, debit jusqu'a 1 Mbit/s.\\
Structure de trame : SOF (1) + ID (11) + RTR (1) + IDE+r0 (2) + DLC (4) + Data (0-64) + CRC (15) + ACK (2) + EOF (7) = 111 bits max.

\end{tcolorbox}

% -----------------------------------------------------------------------
%  2. FORMULAIRE SYNTHeTIQUE
% -----------------------------------------------------------------------

\begin{tcolorbox}[
  ficheSection=colElec,
  title={2. Formulaire synthetique},
]
\small
\begin{multicols}{2}
\renewcommand{\arraystretch}{1.45}
\begin{tabular}{@{}p{5cm}l@{}}
\toprule
\textbf{Formule} & \textbf{Usage} \\
\midrule
\multicolumn{2}{@{}l}{\footnotesize\textit{PWM}} \\
$\alpha = t_{on}/T$ & Rapport cyclique \\
$f_{ech} = N_{pos} \times f_{PWM}$ & Resolution angulaire \\
$N_{pos} = (\theta_{max}-\theta_{min})/\delta\theta$ & Nb positions \\
\midrule
\multicolumn{2}{@{}l}{\footnotesize\textit{DSP}} \\
$f_{max} = P \cdot n / 60$ & Frequence codeur \\
Condition : $f_{DSP} \geq f_{max}$ & Choix DSP \\
\bottomrule
\end{tabular}

\columnbreak

\renewcommand{\arraystretch}{1.45}
\begin{tabular}{@{}p{5cm}l@{}}
\toprule
\textbf{Formule} & \textbf{Usage} \\
\midrule
\multicolumn{2}{@{}l}{\footnotesize\textit{API S7-1500}} \\
$T = T_{min} + PIW8/27648 \cdot \Delta T$ & Mise - l'echelle \\
PIW8 $\in [0\,;\,27\,648]$ & Plage module AI \\
\midrule
\multicolumn{2}{@{}l}{\footnotesize\textit{Bus CAN}} \\
Debit max : 1\,Mbit/s & CAN 2.0 \\
Trame max : 111\,bits (8 octets data) & Structure \\
$D_b = 1/T_b$ & Debit depuis chrono. \\
\bottomrule
\end{tabular}
\end{multicols}
\normalsize
\end{tcolorbox}

% -----------------------------------------------------------------------
%  3. QUESTIONS TYPES
% -----------------------------------------------------------------------

\begin{tcolorbox}[
  ficheSection=colElec,
  title={3. Questions types - extraites des sujets reels},
]

% -- QT1 -------------------------------------------------------------

% Mini-procedure de resolution type copie concours
\begin{tcolorbox}[questiontype,title={Procedure type copie concours (5 etapes)}]
\begin{enumerate}[leftmargin=1.5em]
  \item \textbf{Qualifier le modele} : hypotheses, domaine de validite, grandeurs d'entree/sortie.
  \item \textbf{Ecrire les lois} : equations de base avec symboles et unites SI.
  \item \textbf{Isoler l'inconnue} : expression litterale, puis application numerique tracee.
  \item \textbf{Controler l'ordre de grandeur} : coherence physique et unite finale.
  \item \textbf{Conclure concours} : phrase courte orientee decision technique.
\end{enumerate}
\end{tcolorbox}

% Visuel de methode : chemin de resolution CAPET
\begin{center}
\begin{tikzpicture}[font=\small, node distance=7mm and 6mm, >=Stealth]
  \node[draw, rounded corners, fill=colElec!8, inner sep=4pt] (n1) {Notions};
  \node[draw, rounded corners, fill=colElec!8, right=of n1, inner sep=4pt] (n2) {Methode};
  \node[draw, rounded corners, fill=colElec!8, right=of n2, inner sep=4pt] (n3) {Application annale};
  \node[draw, rounded corners, fill=colElec!8, right=of n3, inner sep=4pt] (n4) {Verification jury};
  \draw[->, thick, colElec] (n1) -- (n2);
  \draw[->, thick, colElec] (n2) -- (n3);
  \draw[->, thick, colElec] (n3) -- (n4);
\end{tikzpicture}
\end{center}
\begin{tcolorbox}[questiontype,
  title={QT1 - Registre ADCON0 du PIC - configuration CAN \hfill\normalfont\footnotesize E1/E2 2022 (PIC 18F8723)}
]
\textbf{enonce.}
Pour selectionner l'entree AN2 du PIC 18F8723 et lancer la conversion, remplir le registre ADCON0 champ par champ. Donner la valeur hexadecimale.

\tcbline

\textbf{Corrige.}

Structure ADCON0 : [ADON | GO/DONE | CHS3 | CHS2 | CHS1 | CHS0 | - | ]
\begin{itemize}
  \item ADON = 1 (CAN active)
  \item GO/DONE = 1 (lancer la conversion)
  \item CHS3:CHS0 = 0010 (voie AN2)
\end{itemize}
Registre : $1\,1\,0010\,00_2$ (si les 2 LSB non utilises = 0) = $\mathrm{C8_{16}}$ (verifier le datasheet exact).

Attention : valeur exacte - lire dans la documentation technique fournie avec le sujet.
\end{tcolorbox}

\vspace{4pt}

% -- QT2 -------------------------------------------------------------

\begin{tcolorbox}[questiontype,
  title={QT2 - PWM servomoteur - frequence d'echantillonnage \hfill\normalfont\footnotesize E2 Commune 2025 (Q18-Q19)}
]
\textbf{enonce.}
Signal PWM - 50 Hz pour un servomoteur ($-90^\circ$ - $+90^\circ$). Precision angulaire requise : $\delta\theta = 0{,}1^\circ$. - quelle frequence le recepteur doit-il echantillonner le signal PWM - Le microcontreleur Arduino (84 MHz) est-il compatible - 
\tcbline

\textbf{Corrige.}

Nombre de positions : $N = \dfrac{180^\circ}{0{,}1^\circ} = 1800$.

Frequence d'echantillonnage necessaire :
\[
  f_{ech} = N \times f_{PWM} = 1800 \times 50 = \mathbf{90\,000\,\mathrm{Hz}} = 90\,\mathrm{kHz}
\]

Condition Arduino : $f_{horloge} = 84\,\mathrm{MHz} \gg 90\,\mathrm{kHz}$ -  compatible (si le code de mesure est suffisamment optimise).
\end{tcolorbox}

\vspace{4pt}

% -- QT3 -------------------------------------------------------------

\begin{tcolorbox}[questiontype,
  title={QT3 - Mise - l'echelle automate - programme de conversion \hfill\normalfont\footnotesize E1 2024 (S7-1500)}
]
\textbf{enonce.}
Le module AI du S7-1500 reeoit un signal 4-20\,mA d'un thermocouple type J ($T \in [-20\,;\,+200e]\,\mathrm{C}$). PIW8 $\in [0\,;\,27\,648]$. ecrire la relation de mise - l'echelle.

\tcbline

\textbf{Corrige.}

\[
  T = T_{min} + \frac{PIW8 - 0}{27\,648 - 0}\,(T_{max} - T_{min})
  = -20 + \frac{PIW8}{27\,648} \times 220
\]

En Ladder/Blocs fonctionnels S7-1500 : instruction SCALE ou calcul avec blocs ADD, MUL, DIV avec les coefficients $a = 220/27648$ et $b = -20$.

Verification : $PIW8 = 0$ - $T = -20^\circ$C - ; $PIW8 = 27\,648$ - $T = 200^\circ$C .
\end{tcolorbox}

\end{tcolorbox}

% -----------------------------------------------------------------------
%  4. ATTENTES DU JURY
% -----------------------------------------------------------------------

\begin{tcolorbox}[
  jury,
  title={4. Attentes du jury - Erreurs recurrentes et points de vigilance},
]

\begin{enumerate}[label=\textbf{\arabic*.}, leftmargin=*, itemsep=5pt]

  \item \textbf{Configurer un registre bit par bit.} Justifier chaque bit de ADCON0 (ou tout autre registre de microcontroleur) - ecrire la valeur hex sans justification est insuffisant.

  \item \textbf{Ceblage d'une entree/sortie automate.} Jury 2021 : "Un candidat au CAPET doit etre capable de cebler une entree TOR/analogique d'un API industriel."

  \item \textbf{Confusion API vs microcontroleur.} L'API est robuste, programme en Ladder/Grafcet pour l'industrie. Le microcontroleur est compact, programme en C/Python pour les fonctions embarquees legeres. Ne pas interchanger.

  \item \textbf{Mise - l'echelle API.} Ne pas oublier le decalage (offset) quand la plage ne commence pas - zero (ex. $PIW8 = 0$ - $T = -20^\circ$C). Verifier les deux bornes.

  \item \textbf{Bus CAN vs CAN (convertisseur).} Dans les sujets, les deux coexistent (2026 E1). Distinguer par le contexte : bus de terrain - convertisseur analogique-numerique.

  \item \textbf{FPGA/VHDL.} Absent des sujets reels 2013-2026 mais au programme. Si la question apparaet : entite + architecture, signal, process concurrent vs sequentiel.

\end{enumerate}

\medskip
\begin{tcolorbox}[
  colback=colJuryFrame!12!colJury,
  colframe=colJuryFrame,
  left=4pt, right=4pt, top=2pt, bottom=2pt,
  fontupper=\small\itshape
]
Rapport jury 2023 : "Notions d'algorithmes, variables informatiques, boucles, systemes embarques sont des connaissances de base attendues des candidats."\\
Rapport jury 2021 : "Un candidat au CAPET doit etre capable de cebler une entree ou une sortie TOR/analogique d'un API industriel programmable."
\end{tcolorbox}

\end{tcolorbox}

% -----------------------------------------------------------------------
%  PIED - PReSENCE DANS LES SUJETS
% -----------------------------------------------------------------------

\begin{tcolorbox}[
  enhanced, breakable,
  colback=colGrayLight,
  colframe=colSeparator,
  left=4pt, right=4pt, top=3pt, bottom=3pt,
  before skip=4pt, after skip=0pt,
  fonttitle=\small\bfseries,
  title={Presence dans les sujets analyses (2013-2026)},
]
\begin{center}
\footnotesize\renewcommand{\arraystretch}{1.15}
\begin{tabular}{lccccccccccccc}
\toprule
\textbf{Annee} & 13 & 14 & 15 & 16 & 17 & 18 & 19 & 20 & 21 & 22 & 23 & 24 & 25 \\
\midrule
E1 (commune)   & - & $\checkmark$ & - & - & - & - & - & - & $\checkmark$ & $\checkmark$ & - & $\checkmark$ & - \\
E2 discipl. IE & - & - & - & - & - & - & $\checkmark$ & - & - & - & - & $\checkmark$ & $\checkmark$ \\
E2 appl. comm. & - & - & - & - & - & - & - & - & - & - & - & $\checkmark$ & $\checkmark$ \\
\bottomrule
\end{tabular}
\end{center}
\smallskip
\footnotesize\textbf{Sous-themes frequents :}
registre CAN microcontreleur (PIC ADCON0) - OneWire DS1820 - mise - l'echelle API S7-1500 - PWM servomoteur - DSP frequence - bus CAN trame - LoRaWAN IoT

\medskip
\textbf{Ordres de grandeur (reperes rapides).}
\begin{itemize}[leftmargin=1.3em]
  \item Reseau BT FR : tension phase-neutre typique \SI{230}{\volt}, tension phase-phase typique \SI{400}{\volt}, frequence \SI{50}{\hertz}.
  \item Courants d'etude frequents en sujet : de quelques \si{\ampere} a quelques dizaines de \si{\ampere}.
  \item En conversion d'energie, un rendement attendu reste generalement entre 0{,}8 et 0{,}98 selon la technologie.
\end{itemize}

\medskip
\textbf{Synthese express (a refaire sans support).}
\begin{enumerate}[leftmargin=1.4em]
  \item Poser les hypotheses du modele et les conventions de signe.
  \item Ecrire 2 a 4 relations cles avec unites.
  \item Realiser un calcul numerique complet et verifier l'ordre de grandeur.
  \item Justifier le resultat par un argument technique de niveau CAPET.
\end{enumerate}

\end{tcolorbox}

\CAPETENDFILE









